\subsection{Un criterio para los $\boldsymbol P$-morfismos.}
\label{subsection:IV.7.5-fr}

\begin{env}[7.5.0]
\label{IV.2.s7.5.0-fr}
Este apartado no se utilizará en el resto del capítulo IV y, por tanto,
puede omitirse en una primera lectura. Veremos además más adelante
(7.9.8) que los resultados de este apartado pueden mejorarse
considerablemente cuando se dispone de la «resolución de
singularidades».

En lo que sigue de este apartado consideraremos una propiedad
$\boldsymbol R(A)$ y denotaremos por $\boldsymbol P(Z,k)$ la propiedad
siguiente:
\begin{quote}
«$Z$ es un preesquema localmente noetheriano sobre un cuerpo $k$ y, para
toda extensión \emph{finita} $k'$ de $k$, si se pone $Z'=Z\otimes_k k'$, se
cumple la propiedad $\boldsymbol R(\mathscr O_{Z',z'})$ para todo
$z'\in Z'$.»
\end{quote}
\end{env}
